
\subsection{An explicit sphere sewing dictionary}

\begin{result}{B.28}{Proved}{Finite pants and chamber presentations compared on every generator}
After capping the leg framings, the finite chamber groupoids
of B.18 for $M_{0,4}$ and $M_{0,5}$ have the following
finite pants presentations and exact inverse dictionaries:
\[
\begin{array}{c|r|r|r|r}
 &\text{objects}&F\text{ arrows}&B\text{ arrows}&\text{relators}\\ \hline
 M_{0,4}&12&6&12&7\\
 M_{0,5}&120&120&144&169 .
\end{array}                                                     \tag{B.83}
\]
Here each $B$ is a positive half-turn of a three-boundary
piece, and each $F$ changes one cut in a four-boundary
piece.  The $24$ associativity relators in the second
row are genuine five-point pentagons.  Both comparison
functors and their inverse natural isomorphism are
given below.  The program \texttt{b\_parenthesized.py}
computes every arrow and relation word.

Consequently every based braid word computed in B.21--B.25
has an explicit word in sphere $F$ and $B$ moves by the
displayed substitution.  This proves the sewing dictionary
for those capped based words.  It retains the separate
questions of the original finite ray transports and the
elliptic framed $S$-move coherence.
\end{result}
\begin{proof}
Put $m=n-1$.  The objects are pairs $(p,T)$, where
$p\in\mathfrak S_m$ is an ordering of the finite labelled
points on the real line and $T$ is a full binary
parenthesization of their ordered positions.  Use the
outer boundary as the $n$th leg.  The two shapes for
$m=3$, and the five for $m=4$, give the object counts.

These are concrete pants markings.  Around each
consecutive block in $T$ take a small rounded rectangle
containing its real points; choose the collar width
strictly decreasing with the depth of the block.
Parent rectangles contain child rectangles, and sibling
rectangles are disjoint.  Their boundaries give the
specified pants cuts.  Moving the widths through these
strict inequalities gives an isotopy, so no additional
mapping class is hidden in the choice.

There is a directed $F$ arrow for every occurrence
\[
 ((A\,B)\,C)\longrightarrow(A\,(B\,C)).
\]
It forgets the first internal cut and inserts the second
inside the same sphere with four boundary components.
Its underlying punctured surface map is the identity.
This is the inverse-removal followed by removal convention
already compared in B.14.  For every cherry consisting
of consecutive individual leaves $i,i+1$, include a
positive $B$ arrow interchanging those two points in
their pair disc.  It fixes the outside of that disc.
The endpoint ordering is $ps_i$, and its underlying
braid is exactly $\sigma_i$.

Define the forward functor $\Phi$ to the chamber model
by
\[
 \Phi(p,T)=p,\qquad
 \Phi(F)=1_p,\qquad
 \Phi(B_{p,T,i})=(p,\sigma_i).                    \tag{B.84}
\]
To give its inverse, choose the left-associated shape
$L_m$.  In the finite $F$ graph for each ordering,
choose the breadth-first tree rooted at $L_m$, with
the lexicographic shape and arrow order used in the
program.  Let $q_{p,T}$ be its unique path from
$(p,L_m)$ to $(p,T)$.  For each $i$, let $T_i$ be the
first shape in the program's order having the cherry
$i,i+1$.  The inverse on the chamber generator is
\[
 \Psi(p,\sigma_i)=
 q_{p,T_i}\, B_{p,T_i,i}\,q_{ps_i,T_i}^{-1}
                                                               \tag{B.85}
\]
in traversal order, and $\Psi(p)=(p,L_m)$.
The image of (B.85) under (B.84) is literally $\sigma_i$,
including its source and target.

Here is the complete finite relation list, rather than
an appeal to an unspecified presentation.  First, for
every non-tree $F$ edge $e:(p,T)\to(p,U)$ impose
\[
 q_{p,T}\,e\,q_{p,U}^{-1}=1.
                                                               \tag{B.86}
\]
Second, for every $B$ edge that is not the chosen
representative at $T_i$, impose
\[
 q_{p,T} B_{p,T,i}q_{ps_i,T}^{-1}
   =q_{p,T_i}B_{p,T_i,i}q_{ps_i,T_i}^{-1}.
                                                               \tag{B.87}
\]
Third, substitute (B.85) into every lifted Artin
relation of B.18, and into its one positive full-scalar
relation.  The latter caps the outer boundary as in
B.18.  All these words and their typed endpoints are
listed in \texttt{b\_parenthesized.json}.

Relations (B.86) say exactly that every $F$ path is
its tree path.  When $m=4$, the five-shape $F$ graph
is a pentagon, so there is one such relator for each
of the $24$ label orderings.  Its geometric equality
is B.14's cut-removal pentagon.  Relation (B.87)
compares two actual realizations of the same pair-disc
half-turn; both have the same braid $\sigma_i$.
The remaining Artin relations are the usual disjoint
support and three-strand isotopies.  The program checks
every resulting braid equality by the faithful Artin
action; the last relation maps to the full scalar braid,
whose capping quotient was proved in B.18.

For completeness, contract the $F$ trees.  Relations
(B.86) remove every remaining $F$ generator, and
(B.87) remove all duplicate $B$ generators.  What
remains is exactly B.18's chamber presentation, on
its labelled vertices, arrows and relators.  Thus
$\Psi$ is defined on the presented groupoid.
The composite $\Phi\Psi$ is the identity on every
object and generator.  The paths $q_{p,T}$ give
the natural isomorphism from $\Psi\Phi$ to the
identity: its naturality for $F$ is (B.86), and
for $B$ is (B.87).  This proves the claimed
equivalence with its actual inverse maps.

The finite enumeration gives $6$ $F$ and $12$ $B$
arrows at $m=3$.  There are no relations (B.86) or
(B.87); the six lifted braid relations and the
scalar relation give seven.  At $m=4$ there are
$120$ $F$ and $144$ $B$ arrows, $24$ relations
(B.86), $72$ relations (B.87), $72$ lifted Artin
relations and one scalar relation, totaling $169$.
These counts and all maps were computed and checked.

Finally, take a computed based word from B.21--B.25.
Read its successive adjacent braid crossings with
their changing label order and replace each positive
crossing by (B.85), each negative crossing by the
inverse word with its changed source.  Equation
(B.84) recovers the original word exactly, so this
is the promised explicit sphere sewing factorization.
The non-braid tree transports of the original finite
ray diagrams remain the stated transports of B.25;
no assertion about an unrecorded framing convention
is obtained by dropping them.
\end{proof}
